\section{Related work}
\label{sec:related}

\paragraph{Parameter control and hyper-heuristics} Choosing among operators during a run is the subject of parameter control in evolutionary algorithms~\cite{EibenEtal1999,KarafotiasEtal2015} and of selection hyper-heuristics~\cite{BurkeEtal2013,DrakeEtal2020}, which choose among low-level heuristics online. Adaptive operator selection (AOS)~\cite{FialhoEtal2010} usually combines a credit-assignment rule (the reward of an operator is its recent improvement, sometimes ranked) with a selection rule: probability matching~\cite{Goldberg1990}, adaptive pursuit~\cite{Thierens2005}, ALNS-style weights~\cite{PisingerRopke2007}, bandits~\cite{AuerEtal2002,DaCostaEtal2008}, or reinforcement learning with a state~\cite{BengioBengio2021}. The credit-assignment and selection rules of selection hyper-heuristics are the same family. We include all of these as learners (Section~\ref{sec:learners}); what the protocol adds is the control that each of them is measured against.

\paragraph{Learning inside metaheuristics for scheduling} Q-Learning has been used to select operators inside an IG for the permutation flowshop~\cite{KarimiMamaghanEtal2023} and, closest to this work, the destruction size of an IG for identical machines with setups~\cite{MissaouiOSullivan2026}; both report that the learning-guided method beats its baselines. Where learned neighbourhood selection is reported to beat strong baselines on routing problems, the policy is trained offline across instances~\cite{SonneratEtal2021}, which is the setting of our offline selector (Section~\ref{sec:selector}). Dynamic algorithm configuration (DAC)~\cite{BiedenkappEtal2020,EimerEtal2021} learns control policies across instances with reinforcement learning; the pre-trained Q table of Section~\ref{sec:repairs} is its tabular form, and the offline selector is its static, per-instance counterpart.

\paragraph{Offline configuration, algorithm selection and the oracle} Offline configurators such as irace~\cite{LopezIbanezEtal2016} and ParamILS~\cite{HutterEtal2009} fix the parameters before the run, and tuning and control are two ways of configuring the same algorithm~\cite{EibenSmit2011}; we therefore treat the validated static action and the irace-configured algorithm as the natural competitors of an online learner. Choosing among algorithms per instance is the algorithm selection problem~\cite{Rice1976,KerschkeEtal2019}; its benchmark libraries~\cite{BischlEtal2016} score a selector by the share of the gap between the single best solver and the per-instance oracle (the virtual best solver) that it closes. Our one-seed oracle is that oracle, measured relative to random choice. What is specific here is that the oracle of a \emph{stochastic} algorithm is inflated by run-to-run noise: taking the best of one run per candidate is a selection from noisy estimates, and its expectation exceeds the best expected value (the winner's curse). We treat this as an estimation problem and remove the bias with cross-fitting, an idea used for nuisance estimation in semiparametric inference~\cite{ChernozhukovEtal2018}; Section~\ref{sec:estimator} and Appendix~\ref{app:estimator} give the argument and a simulation.

\paragraph{Benchmarking methodology} The computational-experiment literature has long asked for controlled experiments that isolate the effect of each component, with hypotheses and statistical tests stated beforehand~\cite{Hooker1995,BarrEtal1995}, and for guidelines on how to compare optimisation algorithms~\cite{LaTorreEtal2021,BartzBeielsteinEtal2020}. Reinforcement-learning research has documented how sensitive reported gains are to seeds, implementations and hyperparameters~\cite{Henderson2018}. We follow that spirit: the protocol of Section~\ref{sec:protocol} fixes the controls (random choice at a matched budget, equal CPU, a tuned static competitor), uses equivalence tests with a declared margin~\cite{Lakens2017} so that ``no effect'' is a finding rather than a failure to reject, and records predictions and code hashes before each battery; preregistration separates predictions from post hoc readings~\cite{NosekEtal2018}, and we label every post hoc reading as such.

\paragraph{What is new} Each ingredient of the protocol is established practice in some community. The contributions are their combination for adaptive selectors; the decision/cost decomposition at equal CPU time; the formal treatment and use of a cross-fitted headroom as a necessary-condition diagnostic; and the evidence assembled with it, which we did not find combined in the literature we audited: a comparison of the learner with random choice at equal CPU time, with a validated static action, with a configurator and with published AOS designs, on four hosts.
